\section{Conclusion and Future Work}
\label{sec:conclusion}
We described \hammer{}, an implemented scholarly metadata aggregator that keeps publication-level matching separate from person-level author decisions, and its author-resolution workflow: candidate retrieval from one provider, up to five evidence signals with explicit availability states, a heuristic score that is not a probability, and conservative gates that end in automatic assignment, manual review or an unresolved result. We do not claim a new disambiguation algorithm. We contribute the combination in one provenance-aware workflow and an evaluation design that reports retrieval and selective assignment together, on one frozen candidate pool, with every outcome kept in the denominators and benchmark-derived results kept apart from independently verified ones.

What has been shown is limited and stated with its basis. The pipeline runs end to end on real provider data and replays offline byte for byte. The evaluation tooling behaves as specified on synthetic data and in simulation, and its first real executions exposed two integration defects that were then fixed. On 68 real queries from LAGOS-AND, benchmark-derived and exploratory, we observed three things. The expected profile was missing from the retrieved pool for about a quarter of the queries. When the name, an institution and a year were supplied, automatic assignment was rare (4 of 50 and 0 of 18 queries) because high-scoring candidates were not separated from similarly named competitors. And the heuristic ordering of a pool agreed with the benchmark no more often than the provider's own order. These are indications of where the present design is limited, and not estimates of accuracy.

What has not been shown is the accuracy of the system. The blinded, independent verification that the protocol requires has not been performed, the coverage pilot and the confirmatory sample have not been captured, and the paper therefore contains no verified precision, recall or coverage. The next steps, in order, are: obtain provider access and run the label-free coverage pilot; fix the target number of automatic assignments from reviewer capacity and size the confirmatory sample (about 435 to 950 primary cases for 32 automatic assignments at the observed coverage); capture and replay it under a new freeze; and have independent reviewers verify every automatic assignment and an audit sample. Beyond that, the evidence suggests three directions: replace the fixed weights with scores calibrated on separate development data and test whether they rank better than the provider's order; raise coverage by giving the resolver richer context or resolved institution identifiers, since the competition among same-name candidates is what blocks automation; and extend the evaluation to other benchmarks and providers.
